\documentclass[10pt]{scrartcl}
\usepackage[mdthm,colorsec,noauthor,noasy]{evan}
\usepackage{handout}
\newcommand{\aug}[2]{\left[\begin{array}{@{}#1@{}}#2\end{array}\right]}
\newcommand{\vv}[1]{\mathbf{#1}}
\newcommand{\ip}[2]{\langle #1,#2\rangle}
\DeclareMathOperator{\spn}{Span}\DeclareMathOperator{\Nul}{Nul}\DeclareMathOperator{\Col}{Col}\DeclareMathOperator{\Row}{Row}\DeclareMathOperator{\tr}{tr}\DeclareMathOperator{\proj}{proj}
\title{Math 54 --- Week 10}
\subtitle{Solutions to the problem set}
\author{}\date{}\ihead{}
\begin{document}
\maketitle

\begin{prob}\textbf{\xlink{math54}{2026-10-26-gram-schmidt-least-squares}{prob-P1}{P1}.}\par\idea{Subtract from $\vv x_2$ its projection onto $\vv v_1=\vv x_1$, clear fractions, then normalize (\xlink{math54}{2026-10-26-gram-schmidt-least-squares}{thm-lay-6-11}{Thm.~6.11}).}
$\vv v_1=(1,2,2)$ ($\vv v_1\cdot\vv v_1=9$), $\vv x_2\cdot\vv v_1=8$, so
\[
  \vv v_2=(4,1,1)-\frac89(1,2,2)=\frac79(4,-1,-1).
\] Clearing the fraction leaves $(4,-1,-1)$; normalizing $\|\vv v_1\|=3$ and $\|(4,-1,-1)\|=3\sqrt{2}$ gives the orthonormal basis
\[
  \frac13(1,2,2),\qquad \frac{1}{3\sqrt{2}}(4,-1,-1).
\]
\end{prob}

\begin{prob}\textbf{\xlink{math54}{2026-10-26-gram-schmidt-least-squares}{prob-P2}{P2}.}\par\idea{Gram--Schmidt clearing fractions after each step.}
\begin{steps}
  \ii $\vv v_1=(1,1,1,1)$ ($\vv v_1\cdot\vv v_1=4$).
  \ii $\vv x_2\cdot\vv v_1=3 \implies \vv x_2-\frac34\vv v_1=\frac14(1,1,1,-3)$; take $\vv v_2=(1,1,1,-3)$ ($\vv v_2\cdot\vv v_2=12$).
  \ii $\vv x_3\cdot\vv v_1=2$ and $\vv x_3\cdot\vv v_2=2$, so
    \[
      \vv x_3-\frac24\vv v_1-\frac{2}{12}\vv v_2=\frac13(1,1,-2,0);
    \]
    take $\vv v_3=(1,1,-2,0)$.
\end{steps}
Orthogonal basis: $\{(1,1,1,1),\ (1,1,1,-3),\ (1,1,-2,0)\}$ (dot products pairwise $0$).
\end{prob}

\begin{prob}\textbf{\xlink{math54}{2026-10-26-gram-schmidt-least-squares}{prob-P3}{P3}.}\par\idea{$Q$ comes from P1; then $R=Q^{\mathsf T}A$ (\xlink{math54}{2026-10-26-gram-schmidt-least-squares}{thm-lay-6-12}{Thm.~6.12}).}
With $\vv q_1=\frac13(1,2,2)^{\mathsf T}$ and $\vv q_2=\frac{1}{3\sqrt{2}}(4,-1,-1)^{\mathsf T}$ from P1: $r_{11}=\vv q_1\cdot\vv x_1=3$, $r_{12}=\vv q_1\cdot\vv x_2=\frac83$, and $r_{22}=\vv q_2\cdot\vv x_2=\frac{7\sqrt{2}}{3}$, giving
\[
  Q = \begin{bmatrix}1/3&4/(3\sqrt{2})\\2/3&-1/(3\sqrt{2})\\2/3&-1/(3\sqrt{2})\end{bmatrix},\qquad
  R = \begin{bmatrix}3&8/3\\0&7\sqrt{2}/3\end{bmatrix}.
\]
\end{prob}

\begin{prob}\textbf{\xlink{math54}{2026-10-26-gram-schmidt-least-squares}{prob-P4}{P4}.}\par\idea{Independent columns give unique least-squares via the normal equations (\xlink{math54}{2026-10-26-gram-schmidt-least-squares}{thm-lay-6-13}{Thm.~6.13}).}
With
\[
  A=\begin{bmatrix}1&0\\0&1\\1&1\end{bmatrix},\qquad \vv b=\begin{bmatrix}1\\1\\0\end{bmatrix}:
\]
\[
  A^{\mathsf T}A=\begin{bmatrix}2&1\\1&2\end{bmatrix},\quad A^{\mathsf T}\vv b=\begin{bmatrix}1\\1\end{bmatrix}
  \implies
  \hat{\vv x}=\frac13\begin{bmatrix}2&-1\\-1&2\end{bmatrix}\begin{bmatrix}1\\1\end{bmatrix}=\begin{bmatrix}1/3\\1/3\end{bmatrix}.
\]
Then $\vv b-A\hat{\vv x}=(\frac23,\frac23,-\frac23)^{\mathsf T}$, giving least-squares error $\|\vv b-A\hat{\vv x}\|=\frac{2\sqrt{3}}{3}$.
\end{prob}

\begin{prob}\textbf{\xlink{math54}{2026-10-26-gram-schmidt-least-squares}{prob-P5}{P5}.}\par\idea{Solve $R\hat{\vv x}=Q^{\mathsf T}\vv b$ by back-substitution (\xlink{math54}{2026-10-26-gram-schmidt-least-squares}{thm-lay-6-15}{Thm.~6.15}).}
For $\vv b=(1,0,0)$, $Q^{\mathsf T}\vv b=(\frac13,\ \frac{4}{3\sqrt{2}})^{\mathsf T}$, which is $(\frac13,\ \frac{2\sqrt{2}}{3})^{\mathsf T}$. From $R\hat{\vv x}=Q^{\mathsf T}\vv b$:
the second row gives $\tfrac{7\sqrt{2}}{3}\hat{x}_2=\tfrac{2\sqrt{2}}{3}$, so $\hat{x}_2=\frac27$; then the first row gives $3\hat{x}_1+\frac83\cdot\frac27=\frac13$, so $\hat{x}_1=-\frac17$.
Check normal equations: $A^{\mathsf T}A=\begin{psmallmatrix}9&8\\8&18\end{psmallmatrix}$, $A^{\mathsf T}\vv b=\begin{psmallmatrix}1\\4\end{psmallmatrix}$, and $A^{\mathsf T}A(-\frac17,\frac27)^{\mathsf T}=(1,4)^{\mathsf T}$.
\end{prob}

\begin{prob}\textbf{\xlink{math54}{2026-10-26-gram-schmidt-least-squares}{prob-P6}{P6}.}\par\idea{Orthonormal columns simplify $A^{\mathsf T}A=I$ immediately.}
Since the columns of $A$ are orthonormal, $A^{\mathsf T}A=I$ (W4) and columns are independent (W3), so the normal equations become $I\vv x=A^{\mathsf T}\vv b \implies \hat{\vv x}=A^{\mathsf T}\vv b$, and $\proj_{\Col A}\vv b=A\hat{\vv x}=AA^{\mathsf T}\vv b$.
\end{prob}

\begin{prob}\textbf{\xlink{math54}{2026-10-26-gram-schmidt-least-squares}{prob-P7}{P7}.}\par\idea{Convert polynomials to value vectors at $0,1,2$ and take standard dot products.}
\begin{steps}
  \ii $p=1+t \mapsto (1,2,3)$ and $q=t^2 \mapsto (0,1,4)$, so $\ip{p}{q}=0+2+12=14$, $\|p\|=\sqrt{14}$, $\|q\|=\sqrt{17}$.
  \ii Gram--Schmidt: $p_0=1\mapsto(1,1,1)$ ($\ip{1}{1}=3$); $\ip{t}{1}=3 \implies p_1=t-1\mapsto(-1,0,1)$ ($\ip{p_1}{p_1}=2$). Finally $\ip{t^2}{1}=5$ and $\ip{t^2}{t-1}=4$, so
  \[
    p_2 = t^2 - \frac53\cdot 1 - \frac42(t-1) = t^2-2t+\frac13.
  \]
  Check $p_2\mapsto(\frac13,-\frac23,\frac13)$, dot~$0$ with $(1,1,1)$ and $(-1,0,1)$.
\end{steps}
\end{prob}

\begin{prob}\textbf{\xlink{math54}{2026-10-26-gram-schmidt-least-squares}{prob-P8}{P8}.}\par\idea{Apply Cauchy--Schwarz (\xlink{math54}{2026-10-26-gram-schmidt-least-squares}{thm-lay-6-16}{Thm.~6.16}) with the constant $1$ vector/function.}
\begin{steps}
  \ii \textbf{(a)} In $\RR^3$, apply Cauchy--Schwarz to $\vv u=(1,1,1)$ and $\vv v=(a,b,c)$: $|a+b+c|\le \sqrt{3}\sqrt{a^2+b^2+c^2}$; squaring gives $(a+b+c)^2\le 3(a^2+b^2+c^2)$.
  \ii \textbf{(b)} In $C[0,1]$ with $\ip{f}{g}=\int_0^1 f(t)g(t)\,dt$, take $g(t)=1$ ($\|1\|=1$):
    \[
      \left|\int_0^1 f(t)\,dt\right|\le 1\cdot\left(\int_0^1 f(t)^2\,dt\right)^{1/2};
    \]
    squaring gives the inequality.
\end{steps}
\end{prob}

\end{document}
